\section{Temporal sampling and observation freshness}
\label{sec:C05:time}
Video and streaming interfaces require a choice of what to observe and when. Uniform frame sampling can miss short events; dense sampling can exhaust computation and context. For an event of duration $d$ observed by periodic samples spaced $\Delta$ apart, a randomly phased event has detection probability $\min(1,d/\Delta)$ under ideal instantaneous sampling and visibility assumptions. This is a calculated sampling fact, not a benchmark result. It shows why an agent that must notice transient confirmations or failures needs event-driven observations or a sampling rate matched to the event timescale.

Freshness is distinct from recognition accuracy. Let an observation be captured at time $t_o$ and the action execute at $t_a$. Its age is $t_a-t_o$. A correct recognition at $t_o$ can become wrong at $t_a$ when another actor changes the environment. Systems should retain observation versions and use preconditions or refreshed observations before consequential operations. A multimodal model cannot infer an unobserved concurrent change reliably merely because it can describe the previous frame. The control policy should decide when the expected cost of stale evidence exceeds the cost of another observation.

A simple decision rule formalizes that trade. Suppose refreshing costs $c_o$ and reduces expected action loss from $L_{\mathrm{stale}}$ to $L_{\mathrm{fresh}}$. Refreshing is beneficial when $L_{\mathrm{stale}}-L_{\mathrm{fresh}}>c_o$, with all quantities measured in a common utility scale. Estimating the loss reduction is difficult and task-dependent; an arbitrary confidence score is not automatically that estimate. Nevertheless, the formulation separates uncertainty about recognition from uncertainty about environmental change. A high-confidence image classification can coexist with a high need for refresh.

An observation policy can itself be learned, but its training distribution must include the consequences of missed events and unnecessary observations. Rewarding only final task completion may obscure excessive sensing cost. Rewarding fewer frames without a reliability constraint can encourage omission of rare but decisive events. Evaluate sensing decisions with retained raw streams where permissions allow, so that missed information can be distinguished from reasoning failure. For privacy-sensitive environments, an appropriately scoped event record may be preferable to storing all video, but the resulting evidential limits should remain explicit.

